\section{Compatibility authority and the skill supply chain}\label{sec:C15:compatibility}
A skill package can modify how an agent interprets tasks, chooses tools, and handles errors. It is therefore part of the effective policy surface. The package's text may be external data, but loading it into a privileged instruction position can give it operational influence. Compatibility and authority should be checked before activation, rather than inferred from a familiar name or a high retrieval score.

A proposed compatibility relation is $\mathrm{Compat}(s,\theta,h,\mathcal T,A)$, involving the skill, model, harness, tool set, and authorization. The relation can require specific tool schemas, artifact formats, and environmental preconditions. It can also forbid a skill from broadening authority. A procedure that was safe when a tool was read-only can become unsafe if the same name later exposes write effects. Version pinning makes this change detectable but does not prove semantic equivalence across versions.

The current Agent Skills format distinguishes required descriptive metadata from optional compatibility and tool declarations. In particular, its \path{allowed-tools} field is experimental and support varies by agent implementation; structural validation checks the artifact format, not every execution permission~\cite{agentskills2026spec}. A consumer must bind package content and referenced resources to the activated version and declare which fields it actually enforces. A description that says a tool is approved cannot override the dispatcher's current policy. Likewise, loading instructions on demand is a context-state transition under a fixed interpreter; changing executable skill resources changes the governing artifact version.

Progent provides a complementary enforcement principle: tool admissibility is checked by a symbolic policy, with policy expansion distinguished from narrowing~\cite{corpus250411703}. A retrieved skill should not silently expand that allowed action set. Even if the model interprets the skill as an instruction to use a new resource, the monitor can require an authorized scope change. This separates useful procedural guidance from the authority to produce effects.

Persistent procedural attacks exploit exactly this boundary. A malicious or mistaken lesson can survive many sessions and influence tasks whose original user never saw the source. Skill provenance should therefore include the generating trajectory, extractor, feedback source, validation results, and activation history. Repeated retrieval is not independent endorsement. A skill created by another agent is evidence from that agent, not a direct observation by the current one.

Consolidation can spread a defect. If a poisoned skill is merged into a broader procedure, simple deletion of the original node may not remove its influence. Record dependency links through merges and re-extractions. On revocation, invalidate affected descendants and test whether cached prompts or preloaded profiles still carry the procedure. The relevant guarantee is the actual scope of future-use prevention, not an unsupported claim of erasing every trace.

The proposed release gate for a skill includes schema checks, negative applicability cases, authority-preservation tests, representative task validation, and a rollback reference. These checks do not guarantee correctness for all future tasks. They make the acceptance conditions reviewable and reduce the chance that a single unverified episode becomes globally active procedural policy. Chapter~\ref{ch:C28} examines the broader attack and defense setting.

